\documentclass[12pt,a4paper]{article}
\usepackage[utf8]{inputenc}
\usepackage[indonesian]{babel}
\usepackage{amsmath}
\usepackage{amsfonts}
\usepackage{amssymb}
\usepackage{geometry}

\geometry{margin=2.5cm}

\title{Penyelesaian Volume Benda dengan Metode Shell}
\author{Ahmad Fakhrul Bawani}
\date{}

\begin{document}

\maketitle

\section*{Soal}

Diketahui:
\begin{itemize}
  \item Kurva pembatas: $y = 256 - x^2$ dan $y = 3x^2$
  \item Batas wilayah: $x = 0$ dan $x \ge 0$
  \item Sumbu putar: Sumbu-$y$
\end{itemize}
Tentukan volumenya menggunakan metode shell!

\subsection*{1. Menentukan Batas Integrasi}
Batas bawah daerah adalah $x = 0$. Batas atas dicari melalui titik potong antara kedua kurva ($y_1 = y_2$):
\begin{align*}
  256 - x^2 &= 3x^2 \\
  256 &= 4x^2 \\
  x^2 &= 64 \\
  x &= \pm 8
\end{align*}
Karena disyaratkan $x \ge 0$, maka batas atas integrasi yang digunakan adalah $x = 8$. Jadi, daerah tersebut diintegrasikan pada interval $[0, 8]$.

\subsection*{2. Rumus Umum}
Karena diputar mengelilingi sumbu-$y$, fungsi diintegrasikan terhadap $x$:
\begin{equation*}
  V = 2\pi \int_{a}^{b} (\text{jari-jari}) \cdot (\text{tinggi}) \, dx
\end{equation*}
Dimana:
\begin{itemize}
  \item Jari-jari tabung adalah $x$.
  \item Tinggi tabung adalah selisih kurva atas dan kurva bawah: \\
        $\text{tinggi} = (256 - x^2) - 3x^2 = 256 - 4x^2$.
\end{itemize}

\subsection*{3. Solusi Integrasi}
Substitusikan komponen yang diketahui ke dalam rumus:
\begin{align*}
  V &= 2\pi \int_{0}^{8} x \cdot (256 - 4x^2) \, dx \\
  V &= 2\pi \int_{0}^{8} (256x - 4x^3) \, dx \\
  V &= 2\pi \left[ 128x^2 - x^4 \right]_{0}^{8} \\
  V &= 2\pi \left( \left[ 128(8)^2 - (8)^4 \right] - \left[ 128(0)^2 - (0)^4 \right] \right) \\
  V &= 2\pi \left( 128(64) - 4096 \right) \\
  V &= 2\pi (8192 - 4096) \\
  V &= 2\pi (4096) \\
  V &= 8192\pi
\end{align*}

Jadi, volume benda putar tersebut adalah \textbf{$8192\pi$ satuan volume}.

\end{document}